\subsection{Boundary Conditions in Atomistic Model Development}

Boundary conditions are part of the physical model, but the most useful
development order need not be the final production geometry.  A practical
atomistic workflow is to establish the bulk model under periodic boundary
conditions first and then change to open boundaries when the scientific
question concerns a finite sample, an edge, or a spatially nonuniform
process.  This is a verification sequence, not a rule that periodic
boundaries are always more physical~\cite{astro-code,kong2026spatially}.

\paragraph{Periodic bulk-verification stage.}
Periodic boundaries identify opposite sides of the simulation cell, so a
displayed edge is connected to its periodic image and is not a physical
surface.  This makes periodic calculations useful for learning and checking
the lattice construction, bulk coordination, bond multiplicities, pair
counting, translational invariance, energy per cell, and supercell
invariance.  Small deterministic lattices should also verify the orientation
of every wrapped directed bond.  The ASTRO reference implementation and its
field tests treat periodic and nonperiodic neighbor semantics as distinct
cases rather than changing them implicitly~\cite{astro-code}.

\paragraph{Open finite-geometry stage.}
After the structure, interaction topology, conservative fields, dynamical
right-hand side, and integrator have passed their periodic checks, use open
boundaries for spatially resolved finite-device studies.  In an atomistic
open-boundary model, edge spins remain dynamical and a finite-range bond is
omitted when its other endpoint lies outside the simulated geometry.  This
must not be implemented by wrapping the neighbor to the opposite edge,
clamping the edge spin, or inserting a zero-valued ghost spin.  Fixed or
pinned spins describe an additional physical constraint and require their
own justification.  Mixed boundaries are appropriate when the target
geometry is periodic along one direction and finite along
another~\cite{astro-code}.

\paragraph{DMI and physical edge response.}
For the unique-bond DMI convention in
Eq.~\eqref{eq:atomistic-hamiltonian}, every retained directed bond must obey
$\mathbf D_{ji}=-\mathbf D_{ij}$.  An open edge omits DMI bonds that cross
the boundary but does not suppress the response of the remaining edge
spins.  The reduced edge coordination can produce chiral canting,
boundary-localized textures, and edge-assisted nucleation; such effects are
part of the finite magnetic problem rather than automatically numerical
artifacts.  Spatially resolved atomistic SOT calculations with DMI provide
an example in which local switching behavior cannot be inferred safely from
only a spatial average~\cite{kong2026spatially}.  Truncating bulk DMI bonds is
the simplest free-edge approximation; it does not establish that a real
atomic termination has unchanged bulk bond coefficients.

\paragraph{Finite-size and wave checks.}
Open boundaries remove direct periodic wraparound, but two edges can still
interact through a domain wall, soliton, or propagating spin wave when the
sample is too small.  Compare increasing lateral sizes and require both the
central-region trajectory and the selected device observable to converge,
while retaining edge-resolved diagnostics.  Periodic calculations require a
corresponding check that a localized texture does not interact materially
with its periodic images.  Neither an ordinary open boundary nor a periodic
boundary is an absorbing boundary.  If suppression of reflected waves is a
numerical requirement, introduce and validate a separate damping buffer and
report the observable region from which that buffer is
excluded~\cite{zhu2020terahertz,astro-code}.

The recommended sequence is therefore
\begin{equation}
    \text{periodic bulk construction and term checks}
    \;\longrightarrow\;
    \text{open finite-geometry spatial studies},
    \label{eq:atomistic-boundary-workflow}
\end{equation}
with both modes retained as explicit, tested configurations.  A result must
state which mode was used; changing the boundary mode changes the modeled
system and is not merely a solver
option~\cite{astro-code,kong2026spatially}.
